\section{Sources, Provenance and Technology Choices}
\label{sec:fontes}
\subsection{A1 and A2 --- Accommodation and Reviews}
The supplied Inside Airbnb extracts\cite{insideairbnb} correspond to the Lisbon snapshot of 19 September 2026, acquired on 14 October 2026 under CC BY 4.0. They contain 1\,950\,076 reviews and 25\,020 listings. Originals are preserved and identified by SHA-256 hashes.

A1 organises comment text, publication dates and references to listings and people. A2 organises the accommodation catalogue and price-observation families. Contacts, individual review scores, observation timestamps, owner documents and additional price histories are generated fields for the exercise; aggregate ratings and publication dates retain their extract provenance. Their model and assessment are developed in Sections~\ref{subsec:modelos_goncalo} and~\ref{subsec:qualidade_goncalo}.

Elasticsearch provides inverted indexing, BM25 lexical search, temporal filters and highlighting of matching excerpts. Vector indexing complements this access with similarity search. MongoDB accommodates objects and lists for amenities, location and observations, supporting profile, owner and quote queries through explicit references and local aggregations.

\subsection{B1 --- Restaurants and Activity}
The Lisbon sample of \emph{TripAdvisor European restaurants}, published by Stefano Leone on Kaggle\cite{tripadvisorkaggle}, contains 5\,261 establishments. The extract was selected using \texttt{country = Portugal} and \texttt{city = Lisbon}. It represents a May 2021 snapshot, was acquired on 14 October 2026 and has a declared CC0 licence. Cuisine and venue-type labels produce 12\,167 associations. The activity scenario consists of simulated group representatives, reservations and visits.

PostgreSQL represents relationships between establishments, labels and events, with declarative keys, domains and referential integrity. Cross-table checks complement these constraints. The relational structure and functional dependencies are presented in Section~\ref{subsec:modelos_tiago}.

\subsection{B2 --- Cultural Heritage}
The cultural catalogue contains 130 municipal records for Lisbon: 67 museums and 63 monuments. Data comes from the \emph{Museus} and \emph{Monumentos Nacionais} catalogues, respectively the EquipamentosCulturais ArcGIS service (layer 3) and Patrimonio service (layer 12)\cite{dados_abertos_cml,cml_monumentos}. The recorded access date is 14 October 2026, and both catalogues declare CC0. The cultural catalogue was generated on 15 October 2026 and preserves original identifiers, source attributes and acquisition metadata.

BaseX retains the hierarchical XML structure and supports traversal through XPath and XQuery. The project XSD checks structure, types, domains, place keys and the reference to the origin. Optional contacts and source fields retain the representations specified by the schema.
